\chapter{禁食的圣经教导}
\section{旧约}
\subsection{群体禁食}
1. 在每年七月初十日的赎罪日，要刻苦己心，什么工都不可做，守为圣日，这是禁食的日子（利一六：29-34，二三：26-32；耶三六：6）。

　　2. 摩西在西乃山领受神所颁布的律法诫命时，四十昼夜，也不吃饭，也不喝水（出三四：28，二四：18；申九：9，18）。

　　3. 以色列人与便雅悯支派彼此战争，便雅悯支派杀死以色列人拿刀的一万八千人，以色列人上伯特利，在主面前哀哭禁食到晚上（士二○：24-28）。

　　4. 撒母耳带领以色列众人，聚集在米斯巴为离弃偶像，专心归向神，当日禁食认罪（撒上七：2-11）。

　　5. 基列雅比居民为扫罗王的骸骨，安葬在雅比的垂柳树下，禁食七日（撒上三一：7-13）。　

　　6. 当亚哈王，与耶洗别王后要霸占拿八的田园时，也命令耶斯列人宣告禁食（王上二一：5-10）。

　　7. 当约沙法王要抵抗摩押人，亚扪人，米乌尼人的攻击时，在隐基底定意寻求主，在犹大全地宣告禁食，求主帮助他们（代下二○：1-4）。

　　8. 在末底改时代，因为哈曼要杀灭犹大人。所以在波斯境内各省各处的犹大人大大悲哀，禁食，哭泣哀号，穿麻衣（斯四：1-3）。以斯帖在领受违例去见王的事上，也请书珊城所有犹大人，为她禁食祷告三日三夜，不吃不喝。以斯帖自己和宫女，也要这样禁食（斯四：15-16）。当末底改和以斯帖得胜哈曼的阴谋后，他们定亚达月十四，十五两日为普珥节，要禁食呼求（斯九：17-32）。当犹大人被掳到巴比伦的七十年中，他们定四，五，七，十等四个月为禁食的日子（亚七：5；八：19）。但神说，这些禁食悲哀岂是向我禁食么（亚七：5）？如果犹大人真心悔改归向真神，遵行神的命令，那么神就要将禁食的日子，变为犹大家欢喜快乐的日子和欢乐的节期，所以神要求犹大人要喜爱诚实和平（亚八：19）。

\subsection{在旧约个人为某一件事而禁食的例子}

　　1. 大卫为哀悼扫罗王和约拿单，禁食哀号到晚上（撒下一：11-12）。

　　2. 尼希米为犹大省遭大难，受凌辱，耶路撒冷城墻被拆毁，城门被火焚烧而哭泣悲哀几天，在天上的上帝面前禁食祈祷（尼一：2-11）。

　　3. 大利乌王因为误听奸人的阴谋，将但以理下在狮子洞中，他回到宫中，终夜禁食，睡不着觉（但六：16-18）。

　　4. 尼尼微城的王和人民，当听到先知约拿所传神要毁灭该城的信息，他们信服上帝。国王下了宝座，脱下朝服，披麻布，坐在灰中，通告全国人民禁食祷告，各人回头离开所行的恶道，丢弃手中的强暴，希望神转意后悔，不发烈怒。神察看他们实际的悔改，就不将所说的灾祸降与他们（拿三：1-10）。　

　　5. 在先知以赛亚书中，记载有犹大人禁食，却是求利益，勒逼人为他们作苦工。禁食却是互相争竞，以凶恶的拳头打人。这类的禁食是神不理会的（赛五八：3-5）。神所要的禁食是：禁食的人必须使被欺压的得自由，将饼分给饥饿的人，接待飘流的穷人，将衣服给赤身的人遮体，顾恤自己的骨肉至亲，要行公义，发怜悯（赛五八：6-12）。　　
\section{新约对禁食祷告的教导}

　　1. 女先知亚拿不离圣殿，禁食祷告，昼夜事奉神（路二：36-37）。

　　2. 施洗约翰的门徒和法利赛人是常禁食的。但耶稣的门徒却少有禁食的，因此引起法利赛人和文士的批评（可二：18-20）

　　3. 耶稣自己在受试探时，禁食四十昼夜（太四：2；路四：2）。

　　4. 在要赶某一类的鬼，需要祷告禁食（太一七：14-21）。

　　5. 安提阿的教会要差派保罗和巴拿巴二人出去宣教，教会禁食祷告，按手在他们的头上（徒一三：1-3）。

　　6. 保罗和巴拿巴两人，在加拉太省的各教会选立长老，禁食祷告，把他们交托所信的主（徒一四：19-23）。　　

\section{归纳}

　
\subsection{禁食祷告的原因}　
　　1. 因罪要得救赎，必须刻苦己心而禁食。

　　2. 为领受神的命令律法而禁食。就是要领受神的旨意，教训，为了专心而禁食。

　　3. 为救人，离弃偶像，归向真神而禁食。

　　4. 恶人要霸占别人田园，也会命令受欺压人禁食。

　　5. 为抵抗敌人攻击，寻求主的旨意而禁食。

　　6. 为拯救同胞脱离恶人的阴谋而禁食。

　　7. 为哀悼朋友死亡而禁食。

　　8. 为神要降灾的警告，而悔改禁食祷告。

　　9. 为事奉神而禁食。

　　10. 为抵抗撒但的试探而禁食。

　　11. 为同工领受职务而禁食祷告。

\subsection{禁食祷告的态度和原则}
　　1. 是专心，严肃的，所以也称为严肃会（赛一：13；珥一：14，二：15）。

　　2. 如果是悔改罪行，必须立志离弃一切罪行，为罪行哀哭忏悔。从此不再犯罪。所以不是在禁食祷告的时候悔改，更是在往后的日子用不犯罪的行为来证明悔改的心志。

　　3. 为拯救同胞脱离神的审判，必须全体信徒同心禁食祷告一段较长的时日。使同胞受感动悔改归向神。

　　4. 为属灵的战争，要得圣灵的同工，明白圣经真理，运用圣经真理为武器，为得胜的力量和智慧，而不倚靠物质的力量，所以不想去吃喝。

　　5. 要小心，不误用禁食为行恶，得私利的借口。以禁食为敬虔的态度去欺骗人。

　　6. 是向神禁食。而求神的察看。是内心的，不是外表的装作，故意叫人看出禁食的样子（太六：16-18）。

　　7. 找几位同心的信徒为罪人得救，失败的信徒复兴，担子沉重的教牧人员等禁食祷告。这类的禁食可以定期举行。但不宣扬，不夸张，不居功。

\subsection{禁食祷告的地点和时间}
　　1. 最好找安静，不容易受干扰的房间，但要告知家人在什么地方。如果是超过一天的，更要告知家人。

　　2. 有人要去祷告山，这不可非议。但没有祷告山的，只要用心灵和诚实禁食祷告，在那里都一样，因为神是灵，是无所不在的。千万不要以为某一地点特别神圣，在那一地方神会特别垂听。

　　3. 时间长短可随圣灵的引导，和内心实际情形的需要而决定；先禁食一，二天，决定再延长。千万不可与别人比赛禁食时间的长短。

\subsection{禁食祷告当注意的事}
　　1. 试探会更大，更严苛。禁食后的失败会更惨。

　　2. 不可抄袭别人禁食的方法，或方式，凡事顺个人的情形而决定。没有一个禁食的完全方式。

　　3. 禁食中得到的真理要自己切实施行，然后分享给别人共同实行。

　　4. 寻求主和祂的旨意为最高的目标。

　　5. 以荣耀主，高举主为最终的目的和动机。